\chapter{Reinforcement learning}
\label{ch:rl}

\begin{comment}
    What is reinforcement learning in words, no equation or algorithms.
    
    How it is different from other types of machine learning.
     - Supervised learning
     - Unsupervised learning
    
    The different elements in reinforcement learning.
     - markov decision process
     - exploration
     - exploitation
     - agent   
     - environment
      • state
      • action
      • reward
     - policy
      • e-greedy
     - value
     - reward
      • different approaches to define rewards
        · continuos rewards
        · reward only at goal
        · negative reward while still alive
     - optional model
     - on/off policy
     - etc

    Very simple example of reinforcement learning.
     - tic-tac-toe
     - 1D travel (1D grid world, wins if it reaches one end, dies if it reaches the other.)
     - quad grid nav (this is more complex that what I would like to have here)
\end{comment}


% Introduction
So far the methods discussed for solving the orbital ascent optimal control problem only yield open-loop solutions. These are solutions for a single trajectory defined by an initial state and the controls as a function of time. In contrast methods such as \acrfull{dp} and \acrfull{rl} produce closed-loop solutions consistent of mappings known as policies that map the controls, also known as actions, with the set of possible states instead of time.

\acrshort{dp} was originally developed in the 1950's as a set of methods used to solve optimal control problems and are widely considered the only feasible way to solve a class of discrete stochastic optimal control problems known as \acrfullpl{mdp}. However \acrshort{dp} suffers from what its author called ``\emph{the curse of dimensionality}'', meaning the computational requirements grow exponentially with the number of state variables. The main reasons are because \acrshort{dp} methods require tables large enough to hold values for the full state set and also require repeated sweeps of the entire state set. 

\acrshort{rl} extends \acrshort{dp} by adding the ability to learn directly from raw experience. Learners use this ability to learn policies by trying out different actions and discovering the most valuable. It does this by not only looking at the reward it received by its latest action, but it also takes into account any future rewards it could receive by taking that action. This learning method coupled with the use of function approximation methods such as \acrfullpl{ann} allows learners to learn useful policies without visiting the entire state set by generalizing between visited states and actions.

This opens up the possibility to compute closed-loop solutions for the high-dimensional orbital ascent trajectory problem that are able to represent a range of optimal or close to optimal trajectories for a wide range of initial states.

This chapter starts with a literature review of the current state-of-the-art developments in \acrshort{rl} and continues with detailed descrip

The first few sections cover the basic theories behind \acrlong{rl}, followed by a description of the more advanced methods to be used during this thesis.
